\subsection{Limitations and Future Work}

Several limitations define the interpretation of this study. First, the
dataset contains 19 laboratory observations, which support common-survival and
observation-boundary claims but not a universal mechanism claim. These
observations are separate experimental recordings rather than independent
biological populations. Second, T1 is a constructed residual depending on
neighborhood selection, finite lag, local affine fitting, and tangential
projection. It is not a directly observed force, intention, or individual
decision variable. Third, the analysis is observational and cannot establish
individual-level causality.

The negative tests are also definition-bound. Failure of the tested
$(C,\dot{C},R)$ closure does not rule out higher-dimensional, delayed,
network-based, or observation-specific stochastic descriptions. Failure of the
state-matched near-pre test does not imply that transitions have no special
dynamics. Propagation was not tested in a confirmatory route and therefore
remains outside the current supported claim.

The submission-hardening checks sharpened these limitations rather than
removing them. The completed full-pipeline pseudo-event calibration supported
the survival result: the observed both-scale count of 14 exceeded the null
maximum of 12 across $B=1000$ replicates; equivalently, no null replicate
reached 14, giving a plus-one empirical $p_{\mathrm{emp}}\approx0.001$. This
strengthens the observation-level null calibration, but it
does not remove the other boundaries. The detrending challenge showed that
the signal was not erased, but the exact both-scale count depended on the
preprocessing choice: centered detrending gave $14/19$, past-only detrending
gave $11/19$, and no rolling detrending gave $13/19$. The local affine
conditioning quality control was reassuring, with all 38 observation-scale
combinations passing and no sampled condition number above 100, but this is a
numerical QC result rather than a biological mechanism result.

Future work should proceed along two broad directions. First, comparative
applications should test whether local affine subtraction reveals similar
residual layers in other weakly ordered collectives. Second, perturbation
experiments or richer measurements of environmental and interaction variables
are needed to determine what biological processes generate T1. Observation
stratification within this dataset can guide both directions by identifying
robust survivors, fragile cases, and stable failures.
